% Problem record op_64db785888803dd3. This TeX file is derived from
% database/problems_json/op_64db785888803dd3.json by scripts/sync-tex.mjs; edit the JSON record
% and rerun the script rather than editing this file.
\section{One-bit simulation of partially entangled qubits}
\label{sec:op_64db785888803dd3}

\paragraph{Problem.}
Can shared randomness and one classical bit exactly simulate every pair of
local projective measurements on every pure entangled two-qubit state?  In
Schmidt form the state is
\begin{equation}
  \lvert\psi_p\rangle
  =\sqrt p\,\lvert00\rangle+\sqrt{1-p}\,\lvert11\rangle,
  \qquad \frac12<p<1.
  \label{eq:p27-schmidt-state}
\end{equation}
For arbitrary Bloch vectors $\mathbf x,\mathbf y\in S^2$ and outcomes
$a,b\in\{0,1\}$, the target distribution associated with
Eq.~\eqref{eq:p27-schmidt-state} is
\begin{equation}
  \begin{aligned}
  P_p(a,b\mid\mathbf x,\mathbf y)
    &=\operatorname{Tr}\!\left[
      \lvert\psi_p\rangle\!\langle\psi_p\rvert
      \bigl(M_{a\mid\mathbf x}\otimes M_{b\mid\mathbf y}\bigr)
      \right],\\
  M_{a\mid\mathbf x}
    &=\frac{I+(-1)^a\mathbf x\cdot\boldsymbol\sigma}{2},
  \qquad
  M_{b\mid\mathbf y}
    =\frac{I+(-1)^b\mathbf y\cdot\boldsymbol\sigma}{2}.
  \end{aligned}
  \label{eq:p27-target-correlation}
\end{equation}
The protocol may use unlimited shared randomness; after receiving
$\mathbf x$, Alice sends Bob one bit, and their local outputs must reproduce
Eq.~\eqref{eq:p27-target-correlation} for every pair of measurement
directions.  If no such universal protocol exists, find a finite-setting
linear inequality satisfied by all one-bit protocols and violated by one of
these target distributions.

\subsection*{Status}

Unsolved

\subsection*{Source}

Gisin explicitly asks whether one classical bit simulates every partially
entangled two-qubit state; Renner and Quintino restate the surviving question
after narrowing the parameter range
\sourcecite{ref:p27-gisin}{Gis09},
\sourcecite{ref:p27-renner-quintino}{RQ23}.

\subsection*{Progress}

\begin{itemize}
  \item Gisin posed the separating-inequality question for one-bit-assisted
  correlations \sourcecite{ref:p27-gisin}{Gis09}.  At the maximally entangled
  endpoint $p=1/2$, Toner and Bacon gave an exact one-bit protocol for all
  projective measurements \sourcecite{ref:p27-toner-bacon}{TB03}; this does
  not cover the nonmaximally entangled interval in
  Eq.~\eqref{eq:p27-schmidt-state}.

  \item Renner and Quintino constructed a one-trit protocol for every
  $1/2\leq p\leq1$.  For $1/2<p<1$, their protocol uses only one bit whenever
  \begin{equation}
    \frac{2p(1-p)}{2p-1}\log\!\left(\frac{p}{1-p}\right)+2(1-p)\leq1,
    \label{eq:p27-one-bit-region}
  \end{equation}
  which holds approximately for $0.835\leq p<1$; the product-state endpoint
  $p=1$ requires no communication.  Thus
  Eq.~\eqref{eq:p27-one-bit-region} leaves the interval
  $1/2<p<0.835$ unresolved \sourcecite{ref:p27-renner-quintino}{RQ23}.

  \item Numerical and semianalytical models support one-bit simulation beyond
  the proved range, but do not supply an exact protocol for all directions
  \sourcecite{ref:p27-sidajaya}{SLYS23}.  Conversely, finite-scenario
  inequalities for one-bit-assisted correlations can be enumerated, but no
  known inequality separates the remaining targets in
  Eq.~\eqref{eq:p27-target-correlation}
  \sourcecite{ref:p27-bacon-toner}{BT03}.
\end{itemize}

\subsection*{References}

\begin{enumerate}[label={},leftmargin=4.8em,labelsep=0.5em]
  \footnotesize
  \item[\textup{[Gis09]}]\label{ref:p27-gisin}
  N. Gisin, ``Bell Inequalities: Many Questions, a Few Answers,'' in
  W. C. Myrvold and J. Christian (eds.), \emph{Quantum Reality, Relativistic
  Causality, and Closing the Epistemic Circle}, The Western Ontario Series in
  Philosophy of Science \textbf{73}, 125--138 (Springer, 2009).
  \href{https://doi.org/10.1007/978-1-4020-9107-0_9}{doi:10.1007/978-1-4020-9107-0\_9};
  \href{https://arxiv.org/abs/quant-ph/0702021}{arXiv:quant-ph/0702021}.

  \item[\textup{[TB03]}]\label{ref:p27-toner-bacon}
  B. F. Toner and D. Bacon, ``Communication Cost of Simulating Bell
  Correlations,'' \emph{Physical Review Letters} \textbf{91}, 187904 (2003).
  \href{https://doi.org/10.1103/PhysRevLett.91.187904}{doi:10.1103/PhysRevLett.91.187904};
  \href{https://arxiv.org/abs/quant-ph/0304076}{arXiv:quant-ph/0304076}.

  \item[\textup{[RQ23]}]\label{ref:p27-renner-quintino}
  M. J. Renner and M. T. Quintino, ``The Minimal Communication Cost for
  Simulating Entangled Qubits,'' \emph{Quantum} \textbf{7}, 1149 (2023).
  \href{https://doi.org/10.22331/q-2023-10-24-1149}{doi:10.22331/q-2023-10-24-1149};
  \href{https://arxiv.org/abs/2207.12457}{arXiv:2207.12457}.

  \item[\textup{[SLYS23]}]\label{ref:p27-sidajaya}
  P. Sidajaya, A. D. Lim, B. Yu, and V. Scarani, ``Neural Network Approach to
  the Simulation of Entangled States with One Bit of Communication,''
  \emph{Quantum} \textbf{7}, 1150 (2023).
  \href{https://doi.org/10.22331/q-2023-10-24-1150}{doi:10.22331/q-2023-10-24-1150};
  \href{https://arxiv.org/abs/2305.19935}{arXiv:2305.19935}.

  \item[\textup{[BT03]}]\label{ref:p27-bacon-toner}
  D. Bacon and B. F. Toner, ``Bell Inequalities with Auxiliary
  Communication,'' \emph{Physical Review Letters} \textbf{90}, 157904
  (2003). \href{https://doi.org/10.1103/PhysRevLett.90.157904}{doi:10.1103/PhysRevLett.90.157904};
  \href{https://arxiv.org/abs/quant-ph/0208057}{arXiv:quant-ph/0208057}.
\end{enumerate}

\subsection*{Comment}

The unresolved parameter range is approximately $1/2<p<0.835$.  It is not
known whether one bit always suffices there or whether a finite-setting
inequality can prove that it does not.

\subsection*{Field}

Quantum foundations

\subsection*{Topic}

Bell nonlocality; Quantum communication complexity; Qubit systems

\subsection*{ID}

\texttt{op\_64db785888803dd3}
